% §2.10 expanded — plane-wave spectrum, propagation, far-field link

\section{Angular Spectrum of Radiation}
\label{sec:ch02.10}

Reciprocity fixed the symmetry of source--field coupling; the complementary descriptive question is how directional radiative content is organized on a reference plane before harmonic bases are chosen in Chapter~\ref{ch:03}. Subsections below derive dispersion, the lateral Fourier transform, and the radiative spectrum after $\Pi_{\mathrm{rad}}\mathcal{R}_\omega$. Finite spatial support on $\Omega_s$ limits which $(k_x,k_y)$ sheets can be excited with appreciable amplitude---that limitation feeds the aperture and singular-value discussion in Section~\ref{sec:ch02.7}.

\subsection{Plane-wave decomposition and dispersion}
\label{sec:ch02.10.1}

On a plane $z=z_0$ in a homogeneous region, a trial field $\E(\mathbf{r})=\E_0\,e^{-\jj(k_x x+k_y y+k_z z)}$ satisfies the vacuum Helmholtz equation only if the propagation vector obeys
\begin{equation}
\label{eq:ch2.dispersion-relation}
k_x^2+k_y^2+k_z^2=k_0^2.
\end{equation}
Define the transverse wavenumber $k_t=\sqrt{k_x^2+k_y^2}$ (rad/m). When $k_t\le k_0$, the longitudinal wavenumber
\begin{equation}
\label{eq:ch2.kz-dispersion}
k_z=\pm\sqrt{k_0^2-k_x^2-k_y^2}
\end{equation}
gives real $k_z$ on the propagating sheet: the mode propagates along $z$ with oscillatory phase. When $k_t>k_0$, $k_z$ is purely imaginary and the mode is \emph{evanescent}: amplitude decays as $e^{-|k_z||z-z_0|}$ away from the surface where it is excited \cite{stratton1941,harrington2001}. Evanescent modes carry local spatial detail and reactive energy; they do not contribute to finite outward power through a distant sphere unless regenerated by sources.

Outgoing selection under $\eim$ picks the root of Eq.~\eqref{eq:ch2.kz-dispersion} whose phase corresponds to propagation away from the source-bearing region \cite{stratton1941}. A standing-wave combination of $+k_z$ and $-k_z$ on the propagating sheet is not the isolated retarded field of a compact source at infinity.

\subsection{Lateral Fourier transform on a scan plane}
\label{sec:ch02.10.2}

Near-field measurements record $\E_t(x,y,z_0)$ on a plane. The lateral spectrum is
\begin{equation}
\label{eq:ch2.plane-wave-angular-spectrum-def}
\widetilde{\E}(k_x,k_y;z_0)
=\int_{-\infty}^{\infty}\!\!\int_{-\infty}^{\infty}
\E(x,y,z_0)\,
e^{\jj(k_x x+k_y y)}\,dx\,dy,
\end{equation}
with inverse
\begin{equation}
\label{eq:ch2.plane-wave-angular-spectrum-inverse}
\E(x,y,z_0)
=\frac{1}{(2\pi)^{2}}
\int_{-\infty}^{\infty}\!\!\int_{-\infty}^{\infty}
\widetilde{\E}(k_x,k_y;z_0)\,
e^{-\jj(k_x x+k_y y)}\,dk_x\,dk_y
\end{equation}
under the standard Fourier convention \cite{stratton1941}. Each pair $(k_x,k_y)$ labels a plane-wave channel; Helmholtz propagation along $z$ multiplies a spectral component by
\begin{equation}
\label{eq:ch2.plane-wave-z-propagation}
\widetilde{\E}(k_x,k_y;z)
=
\widetilde{\E}(k_x,k_y;z_0)\,
e^{-\jj k_z(k_t)(z-z_0)},
\end{equation}
with $k_z(k_t)$ from Eq.~\eqref{eq:ch2.kz-dispersion} on the outgoing branch \cite{stratton1941,jackson1999}. On the evanescent sheet $k_t>k_0$, $|k_z|$ is imaginary and the factor in Eq.~\eqref{eq:ch2.plane-wave-z-propagation} decays exponentially with $|z-z_0|$---evanescent content cannot be extrapolated to the far zone without error growth \cite{hansen1988}.

\subsection{Radiative spectrum and far-field amplitude}
\label{sec:ch02.10.3}

Applying $\mathcal{F}_{\mathrm{lat}}$ to $\mathcal{S}_\omega[\Jimp]$ without outgoing radiative selection mixes reactive evanescent sheets with propagating export. The radiative lateral spectrum is
\begin{equation}
\label{eq:ch2.radiative-lateral-spectrum-def}
\widetilde{\E}_{\mathrm{rad}}(k_x,k_y)
=
\mathcal{F}_{\mathrm{lat}}\big[\Pi_{\mathrm{rad}}\,\mathcal{R}_\omega[\Jimp]\big](k_x,k_y),
\end{equation}
where $\mathcal{F}_{\mathrm{lat}}$ denotes the operator in Eq.~\eqref{eq:ch2.plane-wave-angular-spectrum-def}. On the propagating sheet $k_t\le k_0$, $\widetilde{\E}_{\mathrm{rad}}$ connects to the far-zone directional amplitude $a(\hat{\mathbf{k}})$: each outgoing direction $\hat{\mathbf{k}}$ corresponds to a point on the sphere $k_t=k_0\sin\theta$, $k_z=k_0\cos\theta$ in the standard spherical parametrization \cite{stratton1941,jackson1999}. The far-field limit of Section~\ref{sec:ch02.5.2} is the same information expressed in spherical rather than planar coordinates---the phase factor $e^{+\jj k_0\hat{\mathbf{r}}\cdot\mathbf{r}'}$ in the source integral is the Fourier kernel evaluated on the propagating surface.

For evanescent indices,
\begin{equation}
\label{eq:ch2.evanescent-angular-null-law}
\widetilde{\E}_{\mathrm{rad}}(k_x,k_y)=\mathbf{0}
\quad\text{for}\ k_t>k_0,
\end{equation}
by definition of $\Pi_{\mathrm{prop}}$ on homogeneous backgrounds \cite{stratton1941}. A large measured $|\widetilde{\E}(k_x,k_y)|$ for $k_t>k_0$ therefore indicates near-zone reactive content in $\mathcal{S}_\omega[\Jimp]$, not far-zone radiation---a common misinterpretation when scan data are displayed without outgoing selection \cite{hansen1988,devaney2004}.

\begin{figure}[t]
\centering
\ChOneFigBlock{%
\begin{tikzpicture}[ch1fig, node distance=5mm and 7mm]
 \node[ch1box] (prop) {$k_t\le k_0$\\propagating};
 \node[ch1box, right=14mm of prop] (ev) {$k_t>k_0$\\evanescent};
 \draw[ch1flow] (prop.south) -- ++(0,-4mm) -| (ev.south);
 \node[below=2mm of prop,font=\scriptsize] {far-zone accessible};
 \node[below=2mm of ev,font=\scriptsize] {exponential decay};
\end{tikzpicture}%
}
\caption{Partition of the lateral spectrum: propagating sheets carry outgoing energy; evanescent sheets store local spatial information without radiative export.}
\label{fig:ch2.propagating-evanescent-spectrum}
\end{figure}
